\section{Introduction}
\label{sec:intro}

Give a language model the same four demonstrations in a different order, and it may give a different answer.
Few-shot prompting~\cite{brown2020language} is how most practitioners adapt large language models (LLMs) to a new task, but the order of the demonstrations is almost never treated as a variable.
It is fixed by whatever order the examples happened to be listed in.
If order swings accuracy by many points, then reported few-shot numbers carry hidden variance, and practitioners need to know when order search, calibration or order-averaging is worth the effort.

Prior work shows that order can matter a great deal.
\citet{lu2022fantastically} found that some permutations of the same demonstrations give near state-of-the-art accuracy with GPT-2/GPT-3 while others are close to chance.
\citet{zhao2021calibrate} documented a recency bias: GPT-3 tends to repeat the label of the last demonstration.
\citet{liu2024lost} found U-shaped position effects in long contexts.
From these results a common picture has formed: order sensitivity grows with the number of shots up to a point, classification favors the most recent demonstration, generation favors the first, NLI is much more fragile than sentiment, and larger models are more robust.

{\bf Does this picture hold for modern instruction-tuned models, and is it measured correctly?}
We see three gaps.
First, most studies report the best--worst range over the sampled orders. The range grows with the number of draws, and the number of possible orders ($K!$) grows with the number of shots $K$, so trends in the range over $K$ mix model behavior with sampling.
Second, position influence is usually inferred from observational correlations (\eg how often the prediction matches the last label), which are confounded with demonstration content.
Third, few studies ask whether the ``best'' order is a real property of the prompt or a lucky draw on a small test set, and none we know of separate order effects from the numerical non-determinism of batched bf16 inference.

We address these gaps with a single controlled protocol on the \qwen family (0.5B, 1.5B, 3B, 7B, 14B) across sentiment (\sst), two NLI tasks (\rte, \mnli) and dialogue summarization (\samsum).
We test five specific hypotheses (H1--H5, \secref{sec:method}).
We measure sensitivity with an expected pairwise gap (\pgap) that does not depend on the number of sampled orders.
We estimate per-position influence \emph{causally}, by flipping the label of one demonstration or inserting a style marker into one demonstration's summary and measuring the change in the model's output distribution.
We test whether order effects reproduce across halves of the test set, and we re-run identical prompts to measure the numerical noise floor.

Most of the specific hypothesis does not survive.
Order variance is at its plateau already at $K=2$; the ``rise until 4-shot'' shape appears only in the range, and tracks the number of available orders (2, 6, then 8).
NLI is more order-sensitive than sentiment for every model, but the thresholds fail: the mean best--worst range at $K\geq4$ is 5.9--13.3 points on \mnli (not $>15$).
``Last demonstration $\approx$2--3$\times$ the first'' holds in only 2 of 15 model--task cells; 1.5B and 3B models show \emph{anti-recency} on NLI, and summarization shows \emph{recency} (first/last influence ratio 0.64--0.80) rather than primacy.
Larger models are less order-sensitive on \sst and \rte (\pgap falls from 2.10 to 0.39 points on \sst), but not on \mnli, and on \samsum the 14B model is the most order-sensitive.
Finally, for models of 3B and above on \sst, the best and worst orders cannot be told apart from re-shuffled noise (observed/null range ratio 0.89--1.00).

Our contributions are:
\begin{itemize}[leftmargin=*,itemsep=0pt,topsep=0pt]
    \item We conduct a controlled study of demonstration order across 6 shot counts, 4 tasks of 3 types and 5 model sizes from one family, with about 2.36M scored prompts and deterministic decoding.
    \item We show that the commonly reported best--worst range produces a spurious ``rise then plateau'' trend over shot count, and propose using an expected pairwise gap that is comparable across $K$.
    \item We measure per-position influence with causal label-flip and style-marker interventions, and find that the direction of position effects depends on both task and model size, with no universal recency/primacy split.
    \item We complement sensitivity numbers with split-half reliability tests and an identical-prompt noise floor, which show when a ``best order'' is real and when choosing one would overfit.
\end{itemize}

\Secref{sec:related} reviews related work, \secref{sec:method} describes our protocol, \secref{sec:results} reports results for each hypothesis, and \secref{sec:discussion} discusses implications and limitations.
